% Options for packages loaded elsewhere
\PassOptionsToPackage{unicode}{hyperref}
\PassOptionsToPackage{hyphens}{url}
\PassOptionsToPackage{dvipsnames,svgnames,x11names}{xcolor}
\PassOptionsToPackage{space}{xeCJK}
\documentclass[
]{article}
\usepackage{xcolor}
\usepackage[margin=2.5cm]{geometry}
\usepackage{amsmath,amssymb}
\setcounter{secnumdepth}{-\maxdimen} % remove section numbering
\usepackage{iftex}
\ifPDFTeX
  \usepackage[T1]{fontenc}
  \usepackage[utf8]{inputenc}
  \usepackage{textcomp} % provide euro and other symbols
\else % if luatex or xetex
  \usepackage{unicode-math} % this also loads fontspec
  \defaultfontfeatures{Scale=MatchLowercase}
  \defaultfontfeatures[\rmfamily]{Ligatures=TeX,Scale=1}
\fi
\usepackage{lmodern}
\ifPDFTeX\else
  % xetex/luatex font selection
  \setmainfont[]{DejaVu Sans}
  \setmonofont[]{DejaVu Sans Mono}
  \ifXeTeX
    \usepackage{xeCJK}
    \setCJKmainfont[]{Noto Sans CJK SC}
  \fi
  \ifLuaTeX
    \usepackage[]{luatexja-fontspec}
    \setmainjfont[]{Noto Sans CJK SC}
  \fi
\fi
% Use upquote if available, for straight quotes in verbatim environments
\IfFileExists{upquote.sty}{\usepackage{upquote}}{}
\IfFileExists{microtype.sty}{% use microtype if available
  \usepackage[]{microtype}
  \UseMicrotypeSet[protrusion]{basicmath} % disable protrusion for tt fonts
}{}
\makeatletter
\@ifundefined{KOMAClassName}{% if non-KOMA class
  \IfFileExists{parskip.sty}{%
    \usepackage{parskip}
  }{% else
    \setlength{\parindent}{0pt}
    \setlength{\parskip}{6pt plus 2pt minus 1pt}}
}{% if KOMA class
  \KOMAoptions{parskip=half}}
\makeatother
\usepackage{longtable,booktabs,array}
\usepackage{caption}
\captionsetup[table]{skip=6pt}
\newcounter{none} % for unnumbered tables
\usepackage{calc} % for calculating minipage widths
% Correct order of tables after \paragraph or \subparagraph
\usepackage{etoolbox}
\makeatletter
\patchcmd\longtable{\par}{\if@noskipsec\mbox{}\fi\par}{}{}
\makeatother
% Allow footnotes in longtable head/foot
\IfFileExists{footnotehyper.sty}{\usepackage{footnotehyper}}{\usepackage{footnote}}
\makesavenoteenv{longtable}
\setlength{\emergencystretch}{3em} % prevent overfull lines
\providecommand{\tightlist}{%
  \setlength{\itemsep}{0pt}\setlength{\parskip}{0pt}}
\usepackage{bookmark}
\IfFileExists{xurl.sty}{\usepackage{xurl}}{} % add URL line breaks if available
\urlstyle{same}
% fallback for those not using the hyperref driver hyperxmp:
\makeatletter
\@ifundefined{xmpquote}{\newcommand{\xmpquote}[1]{#1}}{}
\makeatother
\hypersetup{
  colorlinks=true,
  linkcolor={Maroon},
  filecolor={Maroon},
  citecolor={Blue},
  urlcolor={Blue},
  pdfcreator={LaTeX via pandoc}}

\author{}
\date{}

\begin{document}

\section{筑基篇课后习题 X：纳维-斯托克斯存在性与光滑性的 pat
重新观测}\label{ux7b51ux57faux7bc7ux8bfeux540eux4e60ux9898-xux7eb3ux7ef4-ux65afux6258ux514bux65afux5b58ux5728ux6027ux4e0eux5149ux6ed1ux6027ux7684-pat-ux91cdux65b0ux89c2ux6d4b}

\begin{quote}
\textbf{声明 (Declaration)}: 本文\textbf{不代表任何数学结论}
(non-mathematical-claim), 仅为基于 Lean 4
形式化的一种\textbf{实验观测报告} (experimental observation report)。

{\def\LTcaptype{none} % do not increment counter
\begin{longtable}[]{@{}
  >{\raggedright\arraybackslash}p{(\linewidth - 6\tabcolsep) * \real{0.2500}}
  >{\raggedright\arraybackslash}p{(\linewidth - 6\tabcolsep) * \real{0.2500}}
  >{\raggedright\arraybackslash}p{(\linewidth - 6\tabcolsep) * \real{0.2500}}
  >{\raggedright\arraybackslash}p{(\linewidth - 6\tabcolsep) * \real{0.2500}}@{}}
\toprule\noalign{}
\begin{minipage}[b]{\linewidth}\raggedright
习题
\end{minipage} & \begin{minipage}[b]{\linewidth}\raggedright
主题
\end{minipage} & \begin{minipage}[b]{\linewidth}\raggedright
Lean 形式化
\end{minipage} & \begin{minipage}[b]{\linewidth}\raggedright
DOI
\end{minipage} \\
\midrule\noalign{}
\endhead
\bottomrule\noalign{}
\endlastfoot
exercise\_i\_10 & 纳维-斯托克斯的 pat 重新观测（R167） &
PatNavierStokesObservation.lean & 10.5281/zenodo.21917243 \\
\end{longtable}
}
\end{quote}

\textbf{Exercise X for the Foundation-Building Chapter: Navier-Stokes
Existence and Smoothness Revisited from the PAT Perspective}

\begin{quote}
习题定位：筑基篇课后习题系列第十题。利用 mechanics-pat-observation skill
观测纳维-斯托克斯存在性与光滑性（Clay 千禧年问题之一）。过程中调试
skill，并研究该问题在物理空间的折叠中丢失了哪些信息、如何找回。全部新定理只锚筑基篇定理（R047/R050/R085/R143/R147/R153/R161/R165）与框架
NS1/NS2（其他会话已证），不用外部引理。
\end{quote}

\emph{2026-08-13 · Internal research exercise · Lean 4 / mathlib v4.32.2
· 6 new theorems PROVED no sorry · 算器神魂论筑基篇课后习题 X }

\begin{center}\rule{0.5\linewidth}{0.5pt}\end{center}

\subsection{习题陈述}\label{ux4e60ux9898ux9648ux8ff0}

\begin{quote}
纳维-斯托克斯存在性与光滑性（Clay 千禧年问题：3D 不可压
NS，初始速度光滑有限能量 ⟹ 全局光滑解存在）：利用 pat
重新观测。唯一论点：\textbf{NS 的折叠丢失 = 时间方向（ν\textgreater0
黏性不可逆）+ 数值互锁（奇点坍缩）；找回机制 = 时间对合 S²=id +
单位圆模长 ‖ρ(g)‖=1。}
\end{quote}

\begin{center}\rule{0.5\linewidth}{0.5pt}\end{center}

\subsection{零、总论：NS 三项 = pat
三个结构}\label{ux96f6ux603bux8bbans-ux4e09ux9879-pat-ux4e09ux4e2aux7ed3ux6784}

NS 方程（3D 不可压）：∂t u + (u·∇)u + ∇p/ρ = νΔu

{\def\LTcaptype{none} % do not increment counter
\begin{longtable}[]{@{}lll@{}}
\toprule\noalign{}
NS 项 & pat 结构 & 锚定 \\
\midrule\noalign{}
\endhead
\bottomrule\noalign{}
\endlastfoot
∂t u（时间导数） & 时间 = 对合对称对 & R147 CausalityTime \\
(u·∇)u（对流） & 锁定方向链（惯性传播） & R050/R153① \\
νΔu（黏性） & 脱离投影（高维耗散） & R161 \\
\end{longtable}
}

\textbf{千禧年问题的 pat 转译}：存在性 =
折叠后仍有解（脱离投影保持互锁）；光滑性 = 数值互锁保持（r·(1/r) = 1
不坍缩）。

\begin{center}\rule{0.5\linewidth}{0.5pt}\end{center}

\subsection{一、折叠丢失 1：时间方向（ν\textgreater0
黏性不可逆）}\label{ux4e00ux6298ux53e0ux4e22ux5931-1ux65f6ux95f4ux65b9ux5411ux3bd0-ux9ecfux6027ux4e0dux53efux9006}

\textbf{★核心：时间反演折叠下，时间导数项翻转而黏性项不翻转------时间方向在黏性中丢失（EulerVsNS
已证）。}

\subsubsection{论证}\label{ux8bbaux8bc1}

\begin{enumerate}
\def\labelenumi{\arabic{enumi}.}
\tightlist
\item
  \textbf{时间反演 = 穿折越}（NS1/NS2，其他会话已证）：时间 t ↦ -t
  是时间轴上的基点穿折越。
\item
  \textbf{时间导数翻转}（NS2
  timeDeriv\_flips\_under\_reversal）：deriv(velRev u) t = -deriv u
  (-t)------时间导数项在反演下翻转。
\item
  \textbf{黏性项不翻转}（NS2
  viscous\_invariance\_of\_reversal）：νΔ(velRev u) = νΔu------Laplacian
  是空间操作，时间反演不变。
\item
  \textbf{⟹ ν=0（Euler）可逆，ν\textgreater0（NS）不可逆}（NS2
  euler\_reversible\_viscous\_irreversible）------\textbf{时间方向在黏性中丢失}（正反时间解不可区分
  → 熵增方向不可逆）。
\item
  \textbf{找回机制：时间对合}（R147 CausalityTime
  time\_dual\_reduces）：-(-t) = t------折叠类 \{t,
  -t\}（R085）中时间方向成对，经对合找回（时间箭头）。
\end{enumerate}

\subsubsection{形式化（PatNavierStokesObservation.lean，1
定理）}\label{ux5f62ux5f0fux5316patnavierstokesobservation.lean1-ux5b9aux7406}

\begin{itemize}
\tightlist
\item
  \texttt{time\_dual\_reduces}：-(-t) = t------时间对合还原（R147）。
\end{itemize}

\textbf{验证}：0 sorry。锚定 CausalityTime.time\_dual\_directions 同型。

\begin{center}\rule{0.5\linewidth}{0.5pt}\end{center}

\subsection{二、折叠丢失
2：数值互锁（奇点坍缩）}\label{ux4e8cux6298ux53e0ux4e22ux5931-2ux6570ux503cux4e92ux9501ux5947ux70b9ux574dux7f29}

\textbf{★核心：光滑性 = 数值互锁保持（r·(1/r) = 1）；奇点 =
数值互锁坍缩（r→0 或 r→∞）。}

\subsubsection{论证}\label{ux8bbaux8bc1-1}

\begin{enumerate}
\def\labelenumi{\arabic{enumi}.}
\tightlist
\item
  \textbf{光滑性 = 数值互锁保持}（R143
  magnitude\_pair\_reduces\_to\_one）：r·(1/r) = 1------数值对称对还原到
  1（能量守恒的 pat 结构）。
\item
  \textbf{奇点 = 数值互锁坍缩}：r = 0（坍缩到折叠类 R085）时 1/r
  无定义；r → ∞（发散）时数值互锁在极限中失去还原------奇点候选 =
  数值互锁坍缩。
\item
  \textbf{存在性 = 折叠后仍有解}（R161
  脱离投影保持）：黏性（耗散）经投影后，剩余互锁仍保持------解在折叠后仍存在。
\end{enumerate}

\subsubsection{形式化（PatNavierStokesObservation.lean，2
定理）}\label{ux5f62ux5f0fux5316patnavierstokesobservation.lean2-ux5b9aux7406}

\begin{itemize}
\tightlist
\item
  \texttt{smoothness\_numeric\_interlock}：r·(1/r) = 1（r ≠
  0）------光滑性 = 数值互锁保持（R143）。
\item
  \texttt{singularity\_fold\_collapse}：r = 0------奇点 =
  数值互锁坍缩到折叠类。
\end{itemize}

\textbf{验证}：0 sorry。锚定 field\_simp / R143。

\begin{center}\rule{0.5\linewidth}{0.5pt}\end{center}

\subsection{三、找回机制：单位圆模长脱离不变}\label{ux4e09ux627eux56deux673aux5236ux5355ux4f4dux5706ux6a21ux957fux8131ux79bbux4e0dux53d8}

\textbf{★核心：单位圆模长 ‖ρ(g)‖ = 1
是折叠后仍可观测的不变量------光滑性的找回 = 相位表示的单位圆模长保持。}

\subsubsection{论证}\label{ux8bbaux8bc1-2}

\begin{enumerate}
\def\labelenumi{\arabic{enumi}.}
\tightlist
\item
  \textbf{单位圆模长}（R165 phaseRep\_on\_circle/R141）：‖exp(iθ)‖ = 1
  对任意相位成立。
\item
  \textbf{脱离不变}（R161）：折叠（时间反演/数值坍缩）后单位圆模长不变------是可观测的找回机制。
\item
  \textbf{与丢失结构的关系}：丢失的信息（时间方向、数值互锁的坍缩极限）不可全部恢复，但单位圆模长（=1）是丢失结构的可观测印记------\textbf{找回
  = 可观测不变量，非丢失信息全部恢复}（诚实边界）。
\end{enumerate}

\subsubsection{形式化（PatNavierStokesObservation.lean，1
定理）}\label{ux5f62ux5f0fux5316patnavierstokesobservation.lean1-ux5b9aux7406-1}

\begin{itemize}
\tightlist
\item
  \texttt{unit\_circle\_recovers}：‖exp(iθ)‖ = 1------找回 =
  单位圆模长（R165）。
\end{itemize}

\textbf{验证}：0 sorry。锚定 norm\_exp\_ofReal\_mul\_I。

\begin{center}\rule{0.5\linewidth}{0.5pt}\end{center}

\subsection{四、全景（组合定理）}\label{ux56dbux5168ux666fux7ec4ux5408ux5b9aux7406}

\textbf{★核心：时间对合还原（R147）∧ 光滑性 = 数值互锁保持（R143）∧ 找回
= 单位圆模长（R165）------NS 折叠丢失（时间方向 +
数值互锁）的找回机制。}

\subsubsection{形式化（PatNavierStokesObservation.lean，1
定理）}\label{ux5f62ux5f0fux5316patnavierstokesobservation.lean1-ux5b9aux7406-2}

\begin{itemize}
\tightlist
\item
  \texttt{ns\_pat\_perspective}：-(-t) = t ∧ r·(1/r) = 1 ∧ ‖exp(iθ)‖ =
  1------NS pat 全景。
\end{itemize}

\textbf{验证}：0 sorry。合取项分别锚 R147/R143/R165。

\begin{center}\rule{0.5\linewidth}{0.5pt}\end{center}

\subsection{五、skill
调试记录（mechanics-pat-observation）}\label{ux4e94skill-ux8c03ux8bd5ux8bb0ux5f55mechanics-pat-observation}

{\def\LTcaptype{none} % do not increment counter
\begin{longtable}[]{@{}
  >{\raggedright\arraybackslash}p{(\linewidth - 4\tabcolsep) * \real{0.3333}}
  >{\raggedright\arraybackslash}p{(\linewidth - 4\tabcolsep) * \real{0.3333}}
  >{\raggedright\arraybackslash}p{(\linewidth - 4\tabcolsep) * \real{0.3333}}@{}}
\toprule\noalign{}
\begin{minipage}[b]{\linewidth}\raggedright
检查项
\end{minipage} & \begin{minipage}[b]{\linewidth}\raggedright
结果
\end{minipage} & \begin{minipage}[b]{\linewidth}\raggedright
调试
\end{minipage} \\
\midrule\noalign{}
\endhead
\bottomrule\noalign{}
\endlastfoot
结构对应 & ✅ NS 三项 = 时间对合/锁定链/脱离投影 & 新增：黏性 =
脱离投影（viscous\_detachment） \\
折叠丢失 & ✅ 时间方向（EulerVsNS 锚定）+ 数值互锁（奇点） &
发现：数值互锁坍缩是光滑性丢失的关键 \\
找回机制 & ✅ 时间对合 + 单位圆模长 & 新增：单位圆模长 = 脱离不变观测 \\
pat 原生 & ✅ 全部 exp 代数，无坐标展开 & 无 \\
诚实边界 & ✅ 千禧年证明标注 CONJECTURE & 无 \\
\end{longtable}
}

\textbf{skill 改进}：新增''折叠丢失两分法''（时间方向丢失 +
数值互锁丢失）------NS
观测显示折叠丢失不单一，需同时分析时间与数值两个层面。

\begin{center}\rule{0.5\linewidth}{0.5pt}\end{center}

\subsection{六、习题解答总结}\label{ux516dux4e60ux9898ux89e3ux7b54ux603bux7ed3}

{\def\LTcaptype{none} % do not increment counter
\begin{longtable}[]{@{}
  >{\raggedright\arraybackslash}p{(\linewidth - 6\tabcolsep) * \real{0.2500}}
  >{\raggedright\arraybackslash}p{(\linewidth - 6\tabcolsep) * \real{0.2500}}
  >{\raggedright\arraybackslash}p{(\linewidth - 6\tabcolsep) * \real{0.2500}}
  >{\raggedright\arraybackslash}p{(\linewidth - 6\tabcolsep) * \real{0.2500}}@{}}
\toprule\noalign{}
\begin{minipage}[b]{\linewidth}\raggedright
论点
\end{minipage} & \begin{minipage}[b]{\linewidth}\raggedright
内容
\end{minipage} & \begin{minipage}[b]{\linewidth}\raggedright
锚定
\end{minipage} & \begin{minipage}[b]{\linewidth}\raggedright
标签
\end{minipage} \\
\midrule\noalign{}
\endhead
\bottomrule\noalign{}
\endlastfoot
NS 三项结构 & ∂t=时间对合 / (u·∇)u=锁定链 / νΔ=脱离投影 & R147/R050/R161
& PROVED \\
折叠丢失：时间 & ν\textgreater0 黏性不可逆（时间导数翻转/黏性不翻转） &
NS2 EulerVsNS & PROVED（其他会话） \\
折叠丢失：数值 & 奇点 = 数值互锁坍缩（r=0/r→∞） & R143/R085 &
OBSERVATION \\
找回：时间 & 时间对合 -(-t)=t（折叠类 \{t,-t\}） & R147 & PROVED \\
找回：数值 & 单位圆模长 ‖ρ(g)‖=1 脱离不变 & R165/R141 & PROVED \\
★全景 & 折叠丢失（时间+数值）∧ 找回（对合+单位圆） & R147/R143/R165 &
PROVED \\
\end{longtable}
}

\textbf{习题的回答（一句话）}：\textbf{NS
在物理空间折叠（时间反演穿折越）中丢失两类信息------时间方向（ν\textgreater0
黏性不可逆，EulerVsNS 已证）与数值互锁（奇点坍缩到折叠类）；找回机制 =
时间对合 S²=id（时间箭头恢复）+ 单位圆模长
‖ρ(g)‖=1（折叠后仍可观测的不变量）。存在性+光滑性的完整证明不在框架能力内（千禧年问题，CONJECTURE）。}

\begin{center}\rule{0.5\linewidth}{0.5pt}\end{center}

\subsection{七、相关对照}\label{ux4e03ux76f8ux5173ux5bf9ux7167}

{\def\LTcaptype{none} % do not increment counter
\begin{longtable}[]{@{}
  >{\raggedright\arraybackslash}p{(\linewidth - 4\tabcolsep) * \real{0.3333}}
  >{\raggedright\arraybackslash}p{(\linewidth - 4\tabcolsep) * \real{0.3333}}
  >{\raggedright\arraybackslash}p{(\linewidth - 4\tabcolsep) * \real{0.3333}}@{}}
\toprule\noalign{}
\begin{minipage}[b]{\linewidth}\raggedright
文献
\end{minipage} & \begin{minipage}[b]{\linewidth}\raggedright
对应内容
\end{minipage} & \begin{minipage}[b]{\linewidth}\raggedright
备注
\end{minipage} \\
\midrule\noalign{}
\endhead
\bottomrule\noalign{}
\endlastfoot
\textbf{Fefferman, C.} 2000. \emph{Existence and Smoothness of the
Navier-Stokes Equation}（Clay 千禧年陈述） & 存在性+光滑性 & pat
观测对象 \\
\textbf{NS1/NS2/NS3}（其他会话，formal/NavierStokes/） &
时间基点/时间反演可逆性 & 本习题锚定 \\
\textbf{R143/R147/R161/R165}（筑基篇） &
数值互锁/时间对合/脱离投影/单位圆模长 & 本框架定理 \\
\end{longtable}
}

\begin{quote}
诚实边界：存在性+光滑性的完整证明（千禧年问题）不在框架能力内------本习题交付结构观测（折叠丢失分析
+ 找回机制），标注 OBSERVATION/CONJECTURE。
\end{quote}

\begin{center}\rule{0.5\linewidth}{0.5pt}\end{center}

\emph{筑基篇课后习题 X · 2026-08-13 · Lean 4 / mathlib v4.32.2 · 6
theorems PROVED no sorry（PatNavierStokesObservation.lean）· 配套 Lean
形式化：formal/Formal/Toolkit/PatNavierStokesObservation.lean（快照见
../lean/PatNavierStokesObservation.lean）}

\end{document}
